\documentclass[fontsize=13pt,a4paper,oneside, DIV=calc, chapterprefix=true]{scrbook} % twoside for draft
 
%\documentclass[13pt,a4paper,oneside]{book} % twoside for draf

%\usepackage{babel}
\usepackage[utf8]{vietnam}
%\usepackage{times}
%\usepackage{graphicx}

\usepackage{husthesis}
\usepackage{enumitem}

 
\usepackage{tipa}
\usepackage{amssymb}
\usepackage{graphicx}
\usepackage{subcaption}
\usepackage{booktabs}
\usepackage{mathptmx}	% same Time New Roma
\usepackage{amsmath}
\usepackage{amsfonts}
%\usepackage{sty/ipa}
%\renewcommand{\rmdefault}{phv} % Arial
%\renewcommand{\sfdefault}{phv} % Arial
%\renewcommand\thesection{\arabic{section}}

\usepackage{array}
\newcolumntype{P}[1]{>{\centering\arraybackslash}p{#1}}
\newcolumntype{M}[1]{>{\centering\arraybackslash}m{#1}}
\usepackage{fancyhdr}
\usepackage{multirow}
\usepackage{algorithm2e}
\usepackage{float}
\usepackage{listings}
\usepackage{hyperref}
\hypersetup{
    colorlinks=true,
    linkcolor=blue,
    filecolor=magenta,      
    urlcolor=cyan,
}


\crname{Tiểu luận khoa học}
\ctname{NGHIÊN CỨU PHÁT TRIỂN \\
MỘT HỆ THỐNG THỬ NGHIỆM}
\cstuname{
Sinh viên 1 - MSV \\
Sinh viên 2 - MSV \\
Sinh viên 3 - MSV \\
Sinh viên 4 - MSV\\
}

\csMajor{Khoa học máy tính và thông tin}
\csProgram{chuẩn}
\csSupervise{TS. Nguyễn Văn A}
\cttime{2024}

\thesislayout

\begin{document}

\coverpage

\secondcoverpage

\frontmatter


\begin{declaration}
Nhóm sinh viên ... cam đoan ...
\end{declaration}

\begin{acknowledgments}
Để hoàn thành  
\end{acknowledgments}

\begin{abstract}
Nội dung chính của tiểu luận ...    
\end{abstract}
 
% Danh sách từ viết tắt
\nomenclature{CIS}{Computer and Information Science}
\printnomenclature

\tableofcontents
% \listoftables
\listoffigures
%\listofalgorithms

\mainmatter

\fancyhead{}  % Clears all page headers and footers
%\rhead{\thepage}  % Sets the right side header to show the page number
%\lhead{}  % Clears the left side page header
%\fancyfoot[positions]{footer}
\renewcommand{\footrulewidth}{0.4pt}

\pagestyle{fancy}  % Finally, use the "fancy" page style to implement the FancyHdr headers
%

\chapter*{Giới thiệu }
\label{chap:intro}
\addcontentsline{toc}{chapter}{\nameref{chap:intro}}
	
{Giới thiệu đề tài}: Đặt vấn đề
 
{Mục tiêu của đề tài}: Giới hạn mục tiêu đề tài

{Cấu trúc báo cáo}
 
\begin{itemize}
\item Giới thiệu 
\item Mục 1: Tổng quan một số công trình nghiên cứu liên quan tới đề tài: hướng tiếp cận, mô hình sử dụng và kết quả đạt được.
\item Mục 2 ...
\item Mục 3 ...
\item Kết luận: Tóm tắt kết quả đạt được, hướng phát triển tương lai
\end{itemize}



\chapter{Tổng quan  }

Chú ý: các tên chương/mục ở đây chỉ là ví dụ minh hoạ cách soạn thảo, không phải tên chương/mục thật. Sinh viên cần đặt tên chương/mục theo cấu trúc mong muốn của mình.

\section{Tổng quan}

Trong phần này ta có thể tham chiếu tới tài liệu kiểu~\cite{szemeredi1983extremal}.




 Báo cáo tối thiểu là 15 trang, được viết bằng latex theo định dạng này.

Cần có một đoạn ngắn giới thiệu về các nội dung trong báo cáo của nhóm.

Ví dụ: Báo cáo gồm có ... phần:
\begin{itemize}
  \item Chương 1: trình bày về ...
  \item Chương 2: trình bày về ...
  \item Chương 3: trình bày về ...
  \item Kết luận
\end{itemize}

\section{Nội dung nghiên cứu 1}
Ví dụ về cách chèn hình ảnh~\ref{hinh1}.


\begin{figure}[ht]
    \centering
    \includegraphics[scale=0.5]{image/bookfull.png}
    \caption{Image example}
\label{hinh1}
\end{figure}

\section{Nội dung nghiên cứu 2}

\subsection{Ví dụ soạn công thức}
Một số ví dụ về các công thức toán học:

$$
l'_1 = \frac{l_1 - \bar{l_1}}{\sigma_{l1}}\quad\quad\quad
l'_2 = \frac{l_2 - \bar{l_2}}{\sigma_{l2}}
$$


Và một công thức khác:

$$
c(l_1, l_2) = -\log\left(1-\int_{l'_1}^{l'_2}e^{\frac{-t^2}{2}}dt\right)
$$

\subsection{Ví dụ soạn công thức 2}
 
Một công thức toán học phức tạp hơn:

$$
m_{ij} = \mathrm{min} \left\{
    \begin{array}{lll}
      m_{i-1,j-1} & + c(a_{i-1}, b_{j-1}) & \textrm{(simple)}\\
      m_{i-1,j}   & + c(a_{i-1}, 0) + pen_{del} & \textrm{(delete)}\\
      m_{i,j-1}   & + c(0, b_{j-1}) + pen_{ins} & \textrm{(insert)}\\
      m_{i-2,j-1} & + c(a_{i-1} + a_{i-2}, b_{j-1}) + pen_{comp} & \textrm{(compression)}\\
      m_{i-1,j-2} & + c(a_{i-1}, b_{j-1} + b_{j-2}) + pen_{exp} & \textrm{(expansion)}\\
      m_{i-2,j-2} & + c(a_{i-1} + a_{i-2}, b_{j-1} + b_{j-2}) + pen_{mix} & \textrm{(mix)}\\
    \end{array}
    \right.
$$

$$\alpha = \frac{1}{\beta} =\sum_{n=0}^{\infty}\frac{1}{n}$$

Bảng~\ref{bang1} minh hoạ cách chèn bảng.

\begin{table}[h!]
\begin{center}
\begin{tabular}{ |c| l| r| c| } % c means center, l for left, r for right
\hline
 & Column 1 & Column 2 & Column 3 \\ 
\hline
 Line 1 & 276.27 & 326.55 & 237.86 \\  
 \hline
\end{tabular}
\end{center}
\caption{Tiêu đề của bảng} 
\label{bang1}
\end{table}

\chapter{Nội dung thực hiện}

\chapter*{Kết luận}
\label{chap:concl}
\addcontentsline{toc}{chapter}{\nameref{chap:concl}}


Trình bày tóm tắt những gì bạn đã làm được

Trình bày về dự định trong tương lai, hướng nghiên cứu tiếp theo.

Tham khảo thêm về cách viết tài liệu sử dụng Latex:

\url{https://www.overleaf.com/learn/latex/Learn_LaTeX_in_30_minutes}

\bibliographystyle{plain} % ieeetr
\bibliography{refs} 

\end{document}